\documentclass[english,twoside, openright, hidelinks, 11 pt, class=report,crop=false]{standalone}
\input{../../preamb}
\input{../bm}

\usepackage{xr}
\externaldocument{geo_opg_bm}

\renewcommand{\opr}[1]{\vspace{20pt}{\color{blue}\textbf{\ref{#1}}}\\}
\renewcommand{\grubr}[1]{\vspace{20pt}{\color{blue}\textbf{Gruble \ref{#1}}}}

\begin{document}

\opr{opggeoviscos}
Se figur under.
\abc{
	\item Av \pyt har vi at
	\[ AC=\sqrt{1^2+1^2}=\sqrt{2} \]
	Dermed er
	\[ \cos 45^\circ = \frac{AB}{AC}=\frac{1}{\sqrt{2}}=\frac{1}{\sqrt{2}}\cdot\frac{\sqrt{2}}{\sqrt{2}}=\frac{\sqrt{2}}{2} \]
	Videre er 
	\[ \sin 45^\circ = \frac{BC}{AC}=\frac{1}{\sqrt{2}}=\frac{\sqrt{2}}{2} \]
	\item Av \pyt har vi at
	\[ DF=\sqrt{1^2+(\sqrt{3})^2}=2 \]
	Dermed er
	\[ \sin 30^\circ=\frac{DE}{DF}=\frac{1}{2} \]
	\item Vi har at
	\[ \cos 30^\circ=\frac{EF}{DF}=\frac{\sqrt{3}}{2} \]
	\item Vi har at
	\[ \sin 60^\circ = \frac{EF}{DF}=\frac{\sqrt{3}}{2} \quad,\quad
	\cos 60^\circ = \frac{DE}{DF}=\frac{1}{2} \]
	\item Finn $\tan 45^\circ$, $\tan 30^\circ$ og $\tan 60^\circ$
	Vi har at
	\begin{align*}
		\tan 45^\circ &= \frac{BC}{AB}=1 \vn
		\tan 30^\circ &= \frac{DE}{EF}=\frac{1}{\sqrt{3}} \\
		\tan 60^\circ &= \frac{EF}{DE}=\frac{\sqrt{3}}{1}=\sqrt{3}
	\end{align*}
} \vsb 
\fig{opggeo1_lf}

\opr{opggeouminv}
Av vinkelsummen trekanten har vi at
\begin{align*}
	180^\circ &= u+v+90^\circ \\
	v &= 90^\circ-u
\end{align*}
Altså er
\[ \cos(90^\circ-u)=\cos v \]
\textsl{Merk: Vi har også at}
\[ \sin v =\frac{AB}{AC}=\cos u \]
\fig{opggeouminv_lf}

\opr{opggeovistan}
\metode{Metode 1}{2cm} \\
Vi har at
\[ \tan v = \frac{BC}{AB} \]
Videre har vi at
\begin{align*}
	\sin v &= \frac{BC}{AC} \\
	BC &= AC\sin v
\end{align*}
Tilsvarende er
\[ AB= AC\cos v \]
Dermed har vi at
\[ \tan v = \frac{AC\sin v}{AC \cos v}=\frac{\sin v}{\cos v} \] \os

\metode{Metode 2}{2cm} \\
Vi har at
\[ \tan v=\frac{BC}{AB}=\frac{\frac{BC}{AC}}{\frac{AB}{AC}}=\frac{\sin v}{\cos v} \]
\fig{opggeovistan_lf}

\opr{opggeoquadrants}
Vis at hvis $u=180^\circ-v$, så er $\sin u=\sin v$ og $\cos u=-\cos v$.

Av \rref{sincostan} og \rref{sincostan2} har vi at
\begin{align*}
	\sin u &=\frac{AC}{OC}=\sin v \vn
	\cos u &= -\frac{AO}{OC}=-\cos v
\end{align*}
\fig{quadrants_lf}

\opr{opggeofindsincostan}
\abc{
	\item Vi har at
	\begin{align*}
		\sin 180^\circ&=\sin(180^\circ-180^\circ)=\sin 0^\circ =0 \vn
		\cos 180^\circ&=-\cos(180^\circ-180^\circ)=-\cos 0^\circ =-1
	\end{align*}
	Da $\tan x = \frac{\sin x}{\cos x}$ (se \refops{opggeovistan}), er
	\begin{align*}
		\tan 180^\circ&=\frac{0}{-1}=0.
	\end{align*}
	\item Vi har at
	\begin{align*}
		\sin 150^\circ&=\sin(180^\circ-150^\circ)=\sin 30^\circ =\frac{1}{2} \vn
		\cos 150^\circ&=-\cos(180^\circ-150^\circ)=-\cos 30^\circ =\frac{\sqrt{3}}{2}\vn
		\tan 150^\circ&= \frac{\frac{1}{2}}{-\frac{\sqrt{3}}{2}}=-\frac{1}{\sqrt{3}}
	\end{align*}
	\item Vi har at
	\begin{align*}
		\sin 135^\circ&=\sin(180^\circ-135^\circ)=\sin 45^\circ =\frac{\sqrt{2}}{2} \vn
		\sin 135^\circ&=-\cos(180^\circ-135^\circ)=\cos 45^\circ =-\frac{\sqrt{2}}{2} \vn
		\tan 135^\circ &= -1
	\end{align*}
	\item Vi har at
	\begin{align*}
		\sin 120^\circ&=\sin(180^\circ-120^\circ)=\sin 60^\circ =\frac{\sqrt{3}}{2} \vn
		\cos 120^\circ&=-\cos(180^\circ-120^\circ)=\cos 60^\circ =-\frac{1}{2} \vn
		\tan 120^\circ&=\frac{\frac{\sqrt{3}}{2}}{-\frac{1}{2}}=-\sqrt{3}
	\end{align*}
}

\opr{1TV23D1opg1}
\abc{
	\item Vi har at
	\[ \sin u = \frac{AB}{BC}=\frac{4}{5}\quad,\quad \cos u = \frac{3}{5} \]
	Dermed er
	Av \pyt\ er
	\begin{align*}
		(\sin u)^2+(\cos u)^2 &=\left(\frac{4}{5}\right)^2+\left(\frac{3}{5}\right)^2 \br
		&=\frac{16}{25}+\frac{9}{25} \br
		&= 1
	\end{align*}
	\item Vi har at 
	\[ \frac{\sin u}{\cos u}=\frac{\frac{4}{5}}{\frac{3}{5}}=\frac{4}{3}=\frac{AC}{AB}=\tan u \]
	\item Vi har at
	\[ \tan v =\frac{AB}{AC}=\frac{3}{4} \]
	Dermed er
	\[ \tan u \tan v = \frac{4}{3}\cdot \frac{3}{4}=1 \]
	\item Vis at likningene fra a), b) og c) gjelder for alle rettvinklede trekanter.
	Av \pyt\ har vi at
	\begin{align*}
		AB^2 + AC^2 &= BC^2 \\
		\frac{AB^2}{BC^2}+\frac{AC^2}{BC^2}&= 1 \br
		\left(\frac{AB}{BC}\right)^2+\left(\frac{AC}{BC}\right)^2 &= 1\br
		(\sin u)^2+(\cos u)^2 &= 1
	\end{align*}
	At $\frac{\sin u}{\cos u}=\tan u$ har vi vist i \refops{opggeovistan}. Videre
	har vi at
	\[ \tan u \tan v = \frac{AC}{AB}\cdot\frac{AB}{AC}=1 \]
}
\begin{figure}
	\centering
	\includegraphics{\figp{1th22opg1_lf}} \qquad
	\includegraphics{\figp{1th22opg1b_lf}}
\end{figure}


\opr{opggeo2r}
\abc{
	\item Vi har at
	\algv{
		AF&=AD=c-r \vn
		FC&=CE=a-r
	}
	Da $ AF+FC=c $, er
	\algv{
		c-r+a-r&=b \\
		c+a-b&=2r
	}
	\item Med $ c $ som grunnlinje har $ \triangle ABC $ høgde $ b $. Av den klassiske arealformelen for en trekant (se \mb) og formelen fra \refop{opggeoabcr} har vi da at
	\alg{
		(a+b+c)r &= ac \\
		r &= \frac{ac}{a+b+c}
	}
	\item Av oppgave a) og b) er
	\algv{
		c+a-b &= \frac{2ac}{a+b+c} \\
		(c+a-b)(a+b+c)&=2ac \\
		(a+c)^2-b^2 &= 2ac \\
		a^2+c^2 &= b^2
	}
	Formelen kjenner vi igjen som Pytagoras' setning.
}

\opr{opggeotanchord}
\fig{geoopgtanchordlf}
Vi setter $ v=\angle BAC $.
Da $ \angle BAC $ er en periferivinkel, er $ \angle BOC=2v $. $ \triangle BCO $ er likebeint, og derfor er $ \angle CBO=90^\circ-u$ (forklar for deg selv hvorfor). Nå har vi at
\[ \angle EBC=90^\circ-\angle CBO=u \]
\newpage
\grubr{opggeoslope}\\
Vi har at
\alg{
	q&=\frac{BC}{AB}=\tan u \vn
	p&=\frac{-EF}{ED}=-\frac{EF}{ED}=-\tan v
}
\fig{opggeoslope_lf}


\grubr{opggeo180minv}
\fig{opggeo180minvlos}
Vi har at $ y=360-x $. Da $ x $ og $ y $ er de tilhørende sentralvinklene til henholdsvis $ v $ og $ u $, er
\alg{
	2u&=360^\circ-2v \\
	y&=180^\circ-v
}

\newpage
\grubr{opggeokordteo}
\fig{opggeokordteolos}
Av \rref{perfvink} har vi at $ \angle CBA=\angle CDA $. Da $ \angle CPA=\angle DPB $ (de er toppvinkler), er dermed $ \triangle PDA \sim \triangle PBC$. Altså har vi at
\[ DP\cdot PC=AP\cdot PB \]


\grubr{opggeovinkar}
\fig{opggeovinkar_lf}
Det blå området kan deles inn i $\triangle AOC$ og en sektor vi skal kalle for sektor $OBC$. Da $\angle BOC$ er sentralvinkelen til periferivinkelen $\angle BAC$, er $\angle BOC=2v$. Dermed er 
\begin{align*}
	A_{\triangle AOC}&=\frac{1}{2} AO\cdot DC\br
	&=\frac{1}{2}r\cdot OC\sin \angle BOC \br
	&= \frac{1}{2}r^2 \sin (2v)
\end{align*}
Videre har vi at (se \mb\ for formelen for arealet til en sektor)
\[ A_{\text{sektor $OBC$}} = \pi r^2 \frac{2v}{360^\circ}=\pi r^2\frac{v}{180^\circ} \]
Dermed er
\[ A_\text{blått område}=r^2\left(\frac{1}{2}\sin(2v)+\frac{\pi v}{180^\circ}\right) \]

\newpage
\grubr{opgeo45vink}
\fig{opggeo45vinklos}
$ \triangle ABE $ er rettvinklet og likebeint. Periferivinklene $ \angle BCE $ og $ \angle BAE $ spenner over samme bue, og dermed er
\nn{
v = \angle BAE=45^\circ
}

\newpage
\grubr{opggeosin18}
\fig{opggeosin18los}
Da $ \angle BAC=\angle CBA=72^\circ $, er $ \triangle ABC $ likebeint ($ AC=BC $) og $ \angle ACB= 36^\circ $. Altså er også $ \triangle BEC $ likebeint ($ EB=EC $). $ \triangle ABC \sim \triangle AEB$ fordi de har $ \angle BAC $ felles, og $ \angle ACB=\angle EBA $. Vi setter $ x=AB $ og $ y=BC $, og får at
\alg{
\frac{AB}{BC}&=\frac{EA}{AB} \br
\frac{x}{y}&=\frac{y-x}{y} \br
xy+x^2-y^2 &=0
}
Av \textit{abc}-formelen har vi at
\alg{
x&=\frac{-y+\pm\sqrt{y^2-4y^2}}{2} \\
&= \frac{-1\pm\sqrt{5}}{2}y
}
Vi forkaster den negative løsningen for $ x $, og får at
\[BD=\frac{x}{2}=\frac{\sqrt{5}-1}{4}y  \]
Da $ \sin 18^\circ=\frac{BD}{BC} $, er
\[ \sin 18^\circ = \frac{\sqrt{5}-1}{4} \]
\newpage
\grubr{opggeoR}
\fig{opggeoRlos}
Av \rref{perfvink} er $ \angle CSA=2\angle CBA $. Da $ \triangle ASC $ er likesidet, er derfor $ \angle DSA=\angle CBA $. Følgelig er $ \triangle ASD\sim \triangle CBE $. Dette betyr at
\alg{
	\frac{a}{EC}&=\frac{r}{\frac{1}{2}b} \\
	r&=\frac{ab}{2EC}
}
Da $ 2A_{\triangle ABC}=EC\cdot c $, er $ EC=\frac{2A_{\triangle ABC}}{c} $, og dermed er
\[ r=\frac{abc}{4A_{\triangle ABC}} \]

\newpage
\grubr{opggeoviscosuv}
\fig{opgcoscos}
I figuren over merker vi oss at
\alg{
EB&=\sin v & AC &= \sin u  \br
OE&=\cos v & OC&=\cos u
}
Da $ \triangle OCA \sim \triangle BEF $, har vi at
\[ FE=\frac{BE}{OC}AC=\frac{\sin v}{\cos u}\sin u \]
Videre har vi at $ EA=OA-OE=1-\cos v $. Tilsvarende er $ CH=1-\cos u $. I tillegg er 
\nn{
DC=FG=(FE+EA)\cos u=\left(\frac{\sin v}{\cos u}\sin u+1-\cos v\right)\cos u
}
Nå har vi at
\alg{
OD &= OH-CH-DC \\
\cos(u+v)&=1-(1-\cos u)-\left(\frac{\sin v}{\cos u}\sin u+1-\cos v\right)\cos u \\
&= \cos u \cos v- \sin u \sin v
}

\newpage
\grubr{opggeotalesTMconv} \\
Gitt $ \triangle ABC$, hvor $ \angle C=90^\circ $, $ a=BC $, $ b=AC $, og $ c=AB $. Da er $ 4A_{\triangle ABC}=2ab $.
Av \grubrs{opggeoR} er da $  r=\frac{c}{2} $. Altså er $ c=2r $, og er dermed en diameter i den omskrevne sirkelen.

\grubr{opggeotangentsq}\\
Av \refops{opggeotanchord} er $\angle ADB=\angle BAC$. Dermed har $\triangle BAC$ og $\triangle BDA$ to parvis like store vinkler (de har $\angle DBA$ felles), og dermed er trekantene formlike. Altså har vi at
\alg{
\frac{AB}{BD}&=\frac{BC}{AB} \br
AB^2 &= BC\cdot BD
}
\fig{opggeotangentsq_los}

\end{document}


